\section{Phân hoạch không gian an toàn}
Thay vì tiếp cận bài toán tránh vật cản theo cách truyền thống là mô tả các ràng buộc không lồi dựa trên bề mặt vật cản, nhóm tác giả đề xuất một phương pháp chuyển đổi không gian cấu hình tự do (collision-free configuration space) thành hợp của một tập hợp hữu hạn các vùng lồi bị chặn (bounded convex regions). Trong mô hình này, các vùng lồi có thể chồng lấn lên nhau , và bài toán hoạch định chuyển động được phát biểu lại dưới dạng một chương trình quy hoạch lồi nguyên hỗn hợp (Mixed-Integer Convex Program - MICP). Cụ thể, các biến nguyên được sử dụng để gán từng đoạn quỹ đạo đa thức (polynomial trajectory segments) vào các vùng lồi an toàn cụ thể này. Phương pháp này mang lại lợi thế lớn về mặt tính toán vì số lượng biến nguyên chỉ phụ thuộc vào số lượng vùng an toàn được tạo ra (sử dụng thuật toán IRIS) thay vì phụ thuộc vào số lượng mặt của vật cản, giúp giải quyết hiệu quả các môi trường phức tạp với hàng trăm vật cản. Hơn nữa, việc ràng buộc quỹ đạo nằm trọn trong các tập lồi đảm bảo an toàn tuyệt đối cho toàn bộ hành trình liên tục, khắc phục nhược điểm "cắt góc" (corner cutting) thường gặp ở các phương pháp chỉ kiểm tra va chạm tại các điểm lấy mẫu rời rạc.